% Intermediate summary matrix for Q5. May later be condensed, moved to an appendix, or redrawn as a figure.
\begin{table}[tbp]
\centering
\footnotesize
\singlespacing
\caption{Working summary for Q5: identity versus reliability. Distinguishing prediction: reliability learning is caller- or context-specific and renews selectively; habituation generalizes across similar callers and recovers with time.}
\label{tab:summary-q5}
\begin{tabular}{@{}p{1.3in}p{2.4in}p{1.4in}p{0.9in}@{}}
\toprule
\textbf{Process} & \textbf{Finding} & \textbf{Where shown} & \textbf{Status} \\
\midrule
Caller identity & Discriminated; renewal when caller changes & Sciurids; rats; prairie voles & Strong \\
\addlinespace[2pt]
Familiarity / kinship & Often no behavioral effect alone & Marmots (familiarity null); ground squirrels (mother/colony null) & Nulls matter \\
\addlinespace[2pt]
Reliability: discount & Responses fall to callers unpaired with predators & Juvenile Richardson's ground squirrels & Field only \\
\addlinespace[2pt]
Reliability: scrutinize & Responses rise to unreliable callers & Yellow-bellied marmots & Field only \\
\addlinespace[2pt]
Direction of reliability & May depend on cost of independent assessment & Prediction; untested & Open \\
\addlinespace[2pt]
Safety updating & Heterospecific chatter speeds relaxation after hawk call & Gray squirrels & Supporting \\
\addlinespace[2pt]
Lab reliability test & Identity present; reliability manipulated as in sciurids & --- & Gap \\
\bottomrule
\end{tabular}
\end{table}
